\chapter[Calibration: prior work and scope]{Calibration, replicated measurements and the limits of extrapolation}
\label{ch:dc-literature}

The calibration extension addresses a specific inference problem. A bounded
observable has the pointwise geometry
\begin{equation}
 \Gamma_f=\gamma(1-f^2),\qquad |f|\leq1,
 \label{eq:dc-literature-structure}
\end{equation}
with an unknown positive scale $\gamma$. Independent detector randomizations
provide repeated readings of the same latent exposure. Constant and timed-step
reference signals constrain a common multiplicative gain and a uniform shutter
width. A finite collection of jointly covered moments then gives an entropy
lower bound, including uncertainty in those calibration quantities. The
structural identity and the supplied generator bounds remain premises of the
physical model. The reference measurements do not establish either premise.

Several ingredients have substantial precedents. The result should be assessed
as their explicit connection to this structural entropy problem, including its
bias bounds, finite-data coverage and failure cases. Neither the use of
replicas, confidence-set projection nor a search of the literature establishes
priority for the complete construction.

\section{What repeated readings identify}

Deconvolution with repeated measurements is established statistical theory.
Delaigle, Hall and Meister~\cite{dcDelaigle2008}, Section~2.1,
equations~(2.1)--(2.4), use repeated additive measurements of one latent
variable to estimate the error characteristic function and then the signal
density. Their particular construction assumes a real, nowhere-zero error
characteristic function; Theorems~3.1--3.3 compare risks with known-error
procedures under smoothness restrictions. More recently,
Capitao-Miniconi, Gassiat and Leh\'ericy~\cite{dcCapitao2025}, Theorem~1,
identify the signal and two independent error laws up to translation under a
signal Laplace-transform condition that includes bounded signals. Their
identification theorem does not require error moments, symmetry or
nonvanishing characteristic functions. Their quantitative estimation results,
including Proposition~2, impose additional class restrictions.

The present acquisition uses a simpler finite-moment consequence of
conditional independence. If each bounded readout $Z_r$ has conditional mean
$A$, the common latent exposure, then
\begin{equation}
 \mathbb E\!\left[\prod_{r=1}^{k}Z_r\,\middle|\,\text{hidden path}\right]
 =A^k
 \label{eq:dc-literature-replica}
\end{equation}
for distinct independent detector randomizations. This identity also covers
conditionally unbiased shot readouts whose distribution depends on $A$;
it does not require an independent additive-noise representation. Boundedness
supplies finite-sample concentration. With unrestricted unknown tails,
unbiased products alone would supply no such uniform numerical guarantee.

These are genuine repeated observations of the same exposure, not independent
systems with the same marginal distribution, and not copies of one stored
number. Shared sensor errors can invalidate the factorization. Repeating an
exposure does not undo its temporal averaging. For the stationary cosine
circle, the one-time marginal is independent of both angular drift and
diffusion scale, so repetitions at one instant cannot identify those dynamics.
The temporal moments and reference signals are additional physical access.

\section{From covered moments to a covered entropy bound}

Inference through moment restrictions remains useful when point identification
fails. Kato, Sasaki and Ura~\cite{dcKato2021}, Theorem~1 and equation~(3),
derive linear restrictions from repeated additive measurements without
requiring the inversion used in conventional deconvolution. Their Theorem~2
gives uniform coverage of at least $1-\alpha-Cn^{-c}$ under specified class
assumptions. This is a quantitative guarantee, although its class-dependent
constants and approximation conditions are not an immediately usable exact
$\alpha$-level guarantee for the present detector.

Here the target statistics form a declared finite family. Bounded-mean
inequalities~\cite{caHoeffding1963,ogMaurer2009} and a joint error allocation
cover the target and reference moments simultaneously. Optimizing over the
resulting nuisance region is legitimate because that region contains the
true quantities on the same event. Sharing records across moments does not
invalidate a union bound; treating a score fitted to those records as though
it had been fixed in advance would require a separate argument. Confidence-set
mapping itself is standard. For example, S\"ohl and Trabs~\cite{dcSohl2016},
equation~(3.6) and Proposition~3.4, transfer an invariant-density confidence
set through the map from density to drift.

The geometric step specific to this construction is the finite polynomial
hierarchy
\begin{equation}
 w_m(f)=\sum_{j=0}^{m}f^{2j},\qquad
 w_m(f)\Gamma_f=\gamma\bigl(1-f^{2m+2}\bigr)\leq\gamma.
 \label{eq:dc-literature-hierarchy}
\end{equation}
This replaces division by the vanishing factor $1-f^2$ with estimable
polynomial moments. At degree $m$, the specified score uses $2m+2$ independent
center readouts. Temporal second moments cover the unknown scale $\gamma$;
replicated higher moments provide the weighted entropy score. The argument
does not substitute a conditional mean diffusivity for a pointwise geometry.
On the cosine-circle model the ideal hierarchy approaches its full entropy
rate. The constructive approach with estimated calibration also increases
replica access, shrinks the time step and shutter width, and increases the
number of trials. It is not a claim of useful sample complexity or consistency
at fixed positive exposure.

\section{Appropriate detector and diffusion baselines}

Berglund~\cite{dgBerglund2010} and Vestergaard, Blainey and
Flyvbjerg~\cite{dgVestergaard2014} derive the effects of localization error
and motion blur on free-Brownian trajectories. Vestergaard et al.,
Section~III~C and Appendix~C, provide covariance estimators and finite-sample
distribution and confidence calculations. These are substantive baselines,
not merely uncorrected mean-square-displacement fits. In their centered
increment model, fixed frame spacing $\Delta$ gives
\begin{equation}
 \begin{split}
 \mathbb E[(\Delta x_n)^2]&=2D\Delta+2q,\qquad
 \mathbb E[\Delta x_n\Delta x_{n+1}]=-q,\\
 q&=\sigma_{\mathrm{loc}}^2-2DR\Delta.
 \end{split}
 \label{eq:dc-literature-brownian}
\end{equation}
Thus diffusion can be estimated without separately identifying the blur
coefficient $R$ and localization variance from these increments. Gaussianity
makes the covariance the complete centered increment law. This confounding
statement is specific to that observation model.

For confinement, Calderon~\cite{dcCalderon2016}, Section~II~A,
equations~(1)--(10), derives a likelihood for a blurred Ornstein--Uhlenbeck
trajectory with localization error, retaining the correlation between
transition noise and exposure noise. The exposure time is supplied; drift,
confinement, diffusion and localization parameters are estimated. Nakakita and
Uchida~\cite{dgNakakita2020}, Theorem~1, instead estimate an unknown
exposure-to-frame ratio in a shrinking-exposure, increasing-duration regime;
their model does not include the general independent localization-noise
channel considered here. Lind\'en et al.~\cite{dcLinden2017}, Methods
equation~(6), use an image model and a Laplace approximation to obtain
pointwise localization uncertainty, with camera-noise calibration supplied.
That method requires pixel information beyond a retained scalar readout.

A free-Brownian covariance estimator is a suitable control for the
free-Brownian acquisition model. It is not an estimator of the same entropy
target for a bounded cosine observation of a circulating circle. A comparison
on that target requires an adapted circle forward model, matched observations
and resources, and uncertainty calculations for its target functional.
Likewise, an Ornstein--Uhlenbeck likelihood is a confined calibration control,
not a hidden-circle entropy estimator without adaptation.

The constant and step references in the present protocol supply independent
metrological information. A common gain $\kappa$ gives mean $\kappa$ on the
unit reference and mean $\kappa\min(1,s/\tau)$ on a unit step read at known
time $s>0$ through a causal uniform shutter of width $\tau$. Inferring
$(\kappa,\tau)$ from these covered means is ordinary calibration. Its role
here is to propagate a joint nuisance region through the entropy certificate.
Common response between reference and target, known reference shape and time
origin, and the uniform shutter model are explicit assumptions; they are not
learned from the target trajectory alone.

\section{Entropy inference and the scope of empirical calibration}

Noise-aware inference of diffusion, currents and entropy already exists.
Frishman and Ronceray~\cite{ogFrishman2020}, equations~(12)--(14) and
Appendices~E--G, treat diffusion and force inference, localization corrections,
projection and uncertainty. Gingrich et al.~\cite{dgGingrich2026}, Section~II,
equations~(13), (15) and (19)--(22), treat noisy thermodynamic inference using
a known memoryless error kernel. Neither comparison warrants claiming that
noise correction or entropy inference starts here. A projected scalar current
can vanish for the reversible cosine record while a bound conditional on the
admissible hidden diffusion geometry is positive: these are different targets.
Known-kernel pair-law recovery also does not establish unrestricted
continuous-time generator identification from a fixed sampling grid.

Conversely, fixed-grid observations are informative in important restricted
classes. Gobet, Hoffmann and Rei\ss~\cite{dcGobet2004}, Section~3 and
Theorems~2.4--2.5, estimate drift and diffusion of a reflected scalar diffusion
on a compact interval using its self-adjoint transition operator and Sobolev
regularity assumptions. Their result prevents a blanket claim that all
fixed-grid diffusion inference is impossible. S\"ohl and
Trabs~\cite{dcSohl2016}, Assumption~3.1, impose drift regularity and recurrence
conditions for their scalar constant-volatility model. Their adaptive bands
add the self-similarity condition in Assumption~4.1; Theorem~4.6 gives
asymptotic coverage and shrinking sup-norm diameter. These are structured
positive results, not permission to replace a required numerical envelope by
an empirical average.

Chapter~\ref{ch:dc-calibration-limits} supplies the complementary failure
case: rare smooth regions can make global diffusivity or generator suprema
arbitrarily large while leaving finite-duration trajectory laws arbitrarily
close. Its confidence consequence uses the standard transfer of probabilities
through total variation~\cite{caLeCamYang2000}. The construction retains
fixed dimension, a common ellipticity lower bound and a common mixing gap;
every individual model is smooth, but the class has no common regularity
envelope. The obstruction concerns unseen global suprema, even with ideal
continuous observations. It neither contradicts pathwise local-volatility
recovery nor identifies the same problem as adaptation of confidence bands
over smoothness classes. All those countermodels are reversible.

The resulting scope is concrete. Supplied geometry and regularity, genuine
replicas, temporal moments and calibrated references can support a computable
finite-data entropy certificate with simultaneous nuisance coverage.
Unsupported global extrapolation cannot. Numerical experiments can test the
implementation and show when a certificate is informative; they do not
validate the physical premises or constitute an empirical discovery. The
primary-source audit records the inspected versions and exact comparison
points without asserting an exhaustive novelty search.
